% !TEX encoding = UTF-8 Unicode
\documentclass[a4paper]{article}

\usepackage{color}
\usepackage{url}
\usepackage[T2A]{fontenc} % enable Cyrillic fonts
\usepackage[utf8]{inputenc} % make weird characters work
\usepackage{graphicx}
\usepackage{listings}
\usepackage{multicol}
\usepackage{geometry}
\usepackage{lipsum}

%\usepackage[nottoc]{tocbibind}

\usepackage[english,serbian]{babel}
%\usepackage[english,serbianc]{babel} %ukljuciti babel sa ovim opcijama, umesto gornjim, ukoliko se koristi cirilica

\usepackage[unicode]{hyperref}
\hypersetup{colorlinks,citecolor=green,filecolor=green,linkcolor=blue,urlcolor=blue}

%\newtheorem{primer}{Пример}[section] %ćirilični primer
\newtheorem{primer}{Primer}[section]

\definecolor{mygreen}{rgb}{0,0.6,0}
\definecolor{mygray}{rgb}{0.5,0.5,0.5}
\definecolor{mymauve}{rgb}{0.58,0,0.82}
\definecolor{greyish}{rgb}{0.97,0.97,0.97}


\lstset{ 
  backgroundcolor=\color{greyish},   % choose the background color; you must add \usepackage{color} or \usepackage{xcolor}; should come as last argument
  basicstyle=\scriptsize\ttfamily,        % the size of the fonts that are used for the code
  breakatwhitespace=false,         % sets if automatic breaks should only happen at whitespace
  breaklines=true,                 % sets automatic line breaking
  captionpos=b,                    % sets the caption-position to bottom
  commentstyle=\color{mygreen},    % comment style
  deletekeywords={...},            % if you want to delete keywords from the given language
  escapeinside={\%*}{*)},          % if you want to add LaTeX within your code
  extendedchars=true,              % lets you use non-ASCII characters; for 8-bits encodings only, does not work with UTF-8
  firstnumber=1,                   % start line enumeration with line 1000
  frame=single,	                   % adds a frame around the code
  keepspaces=true,                 % keeps spaces in text, useful for keeping indentation of code (possibly needs columns=flexible)
  keywordstyle=\color{blue},       % keyword style
  language=Python,                 % the language of the code
  morekeywords={*,...},            % if you want to add more keywords to the set
  numbers=left,                    % where to put the line-numbers; possible values are (none, left, right)
  numbersep=5pt,                   % how far the line-numbers are from the code
  numberstyle=\tiny\color{mygray}, % the style that is used for the line-numbers
  rulecolor=\color{black},         % if not set, the frame-color may be changed on line-breaks within not-black text (e.g. comments (green here))
  showspaces=false,                % show spaces everywhere adding particular underscores; it overrides 'showstringspaces'
  showstringspaces=false,          % underline spaces within strings only
  showtabs=false,                  % show tabs within strings adding particular underscores
  stepnumber=2,                    % the step between two line-numbers. If it's 1, each line will be numbered
  stringstyle=\color{mymauve},     % string literal style
  tabsize=2,	                   % sets default tabsize to 2 spaces
  title=\lstname                   % show the filename of files included with \lstinputlisting; also try caption instead of title
}

\definecolor{dkgreen}{rgb}{0,0.6,0}
\definecolor{gray}{rgb}{0.5,0.5,0.5}
\definecolor{mauve}{rgb}{0.58,0,0.82}


\begin{document}

\title{\textbf{Predviđanje filmskih rejtinga}\\
\small{Projekat u okviru kursa Mašinsko učenje}}

\author{Nenad Perišić\\perisicnenad96@gmail.com\\}

\medskip

\newpage

\maketitle

%\begin{figure}[hb!]
%\begin{center}
%\small{Matematički fakultet \\ u Beogradu\\}
%\includegraphics[scale=0.5]{matf_logo.png}
%\end{center}
%\end{figure}

%\vfill
\vspace{2cm}
\begin{figure}[b!]
\begin{center}
\includegraphics[scale=0.67]{matf_logo.png} \\
\small{Matematički fakultet \\ u Beogradu}
\end{center}
\end{figure}


\pagebreak

\tableofcontents

\newpage

\section{Uvod}
\label{sec:uvod}
neki uvod
\cite{skripta}

\pagebreak

\section{Zaključak}
\label{sec:zakljucak}
U ovom radu smo prikazali neke od tehnika i alata koje se mogu upotrebiti pri klasterovanju ovako 
velikog skupa podataka. Zaključak koji možemo da izvedemo jeste taj da se alat CLUTO pokazao kao najbolje sredstvo za klasterovanje ovakvih vrsta podataka. Ostale metode nisu dale dobre rezulate. Napomenimo i to da smo prilikom klasterovanja ćelija uglavnom dobijali bolje rezultate u odnosu na klasterovanje gena.



\pagebreak

\addcontentsline{toc}{section}{Literatura}
\appendix
\bibliography{reference}
\bibliographystyle{plain}
\appendix

\end{document}
